\section{Introduction}
\subsection{Problem Description}

The proposed system is an RTOS-based smart climate controller that continuously monitors environmental temperature and
humidity using a DHT20 sensor. Based on the measured values, the system automatically controls visual actuator
indicators represented by LEDs for the heater, cooler, and humidifier. The purpose of the system is to maintain
environmental conditions within predefined ranges and demonstrate real-time task scheduling and coordination on an
ESP32-S3 microcontroller.

Temperature and humidity regulation is important in many applications such as agriculture, storage facilities, and
greenhouse management. Maintaining suitable environmental conditions can improve productivity and reduce the risk of
damage caused by excessive heat, low temperature, or insufficient humidity. In this project, the physical actuators are
represented by LEDs; however, the same control logic could be extended to real devices such as fans, heaters, and
humidifiers.

\subsection{Required Features}

The system implements the following required monitoring and control functions:

\begin{itemize}
  \item The DHT20 sensor measures temperature and humidity periodically every 5 seconds.
  \item The measured values are displayed through the serial monitor.
  \item The Heater indicator (LEDs D3 and D4) changes color according to the temperature level:

        \begin{itemize}
          \item GREEN: temperature within the safe operating range.
          \item ORANGE: temperature approaching a warning threshold.
          \item RED: temperature in a critical range.
        \end{itemize}
  \item The Cooler indicator (LEDs D5 and D6) is activated for a fixed duration of 5 seconds when the temperature exceeds the
        cooling threshold. After the activation period, the temperature is measured again to determine whether additional
        cooling is required.
  \item The Humidifier indicator (LEDs D7 and D8) executes a timed sequence when the humidity falls below the specified
        threshold:

        \begin{itemize}
          \item GREEN for 5 seconds
          \item YELLOW for 3 seconds
          \item RED for 2 seconds
        \end{itemize}
  \item After the sequence completes, the humidity value is checked again to determine whether the sequence should be repeated.
  \item The onboard status LED on GPIO48 blinks every second to indicate that the RTOS scheduler is running correctly.
\end{itemize}

\subsection{Additional Proposed Features}

In addition to the required functionality, an intermediate Decision Maker task is introduced between the
sensor-monitoring task and the actuator tasks. The Decision Maker receives the latest temperature and humidity
measurements and determines which actuator tasks should be activated.

This additional component improves synchronization between tasks and prevents unintended repeated activations.
According to the project requirements, sensor monitoring should resume only after the cooler or humidifier has
completed its current operation. If the monitoring task directly signaled actuator tasks, a situation could occur in
which one actuator finishes earlier than another and immediately triggers a new measurement cycle. For example, if the
cooler finishes while the humidifier sequence is still running, the newly measured humidity value could cause the
humidifier task to be activated again before its previous sequence has completed. Such behavior would lead to
overlapping operations and potentially incorrect control decisions.

By introducing the Decision Maker, actuator execution is coordinated more effectively, ensuring that each control
sequence completes before a new activation request is issued. This design demonstrates the use of RTOS synchronization
concepts and improves the reliability of the overall control system.


